\documentclass[UTF8,a4paper,fontset=fandol]{ctexart}
\usepackage[left=16mm,right=16mm,top=15mm,bottom=16mm]{geometry}
\usepackage{fancyhdr}
\usepackage{setspace}
\usepackage{microtype}
\usepackage{hyperref}
% Route circled digits ①–⑳ and star symbols through the CJK font engine
\xeCJKDeclareCharClass{CJK}{"2460 -> "2473, "2605, "2606}
% Fallback to SimSun for rare CJK characters missing from Fandol (declared after ctex loads)
\AtBeginDocument{\setCJKfallbackfamilyfont{\CJKrmdefault}{SimSun}}
\hypersetup{unicode=true,pdftitle={2026年普通高等学校招生全国统一考试（全国I卷）语文},pdfauthor={整理排版学习版}}
\pagestyle{fancy}
\fancyhf{}
\fancyfoot[C]{\small 第\thepage 页}
\renewcommand{\headrulewidth}{0pt}
\renewcommand{\footrulewidth}{0pt}
\setlength{\parindent}{0pt}
\setlength{\parskip}{0pt}
\setlength{\emergencystretch}{2em}
\linespread{1.03}
\begin{document}
\fontsize{8.6pt}{11.2pt}\selectfont
\sloppy
\textbf{绝密★启用前}\par
\vspace{0.35em}\par
\begin{center}\bfseries\fontsize{15pt}{19pt}\selectfont 2026 年普通高等学校招生全国统一考试（全国 I 卷）\end{center}
\fontsize{8.6pt}{11.2pt}\selectfont
\begin{center}\bfseries\fontsize{18pt}{22pt}\selectfont 语文\end{center}
\fontsize{8.6pt}{11.2pt}\selectfont
\begin{center}\bfseries 本试卷共 10 页，22 小题，满分 150 分。考试用时 150 分钟。\end{center}
\vspace{0.35em}\par
整理说明：本文件为可检索的重新排版学习版，并非教育部教育考试院发布的原版扫描件。题目文本依据公开完整转录整理，并结\par
合教育部教育考试院相关命题评析与公开题目材料交叉核验。原题中用于第18题辨析的四处错别字予以保留。\par
\vspace{0.35em}\par
\vspace{0.35em}\par
\textbf{注意事项：}\par
1．答卷前，考生务必用黑色字迹的钢笔或签字笔将自己的姓名、考生号、考场号和座位号填写在答题卡上。\par
2．作答选择题时，选出每小题答案后，用 2B 铅笔把答题卡上对应题目选项的答案信息点涂黑；如需改动，用橡皮擦\par
干净后，再选涂其他答案，答案不能答在试卷上。\par
3．非选择题必须用黑色字迹的钢笔或签字笔作答，答案必须写在答题卡各题目指定区域内相应位置上；如需改动，\par
先划掉原来的答案，然后再写上新的答案；不准使用铅笔和涂改液。不按以上要求作答的答案无效。\par
4．考生必须保持答题卡的整洁。考试结束后，将试卷和答题卡一并交回。\par
\vspace{0.35em}\par
\vspace{0.35em}\par
\vspace{0.25em}\textbf{一、阅读（74 分）}\par
\vspace{0.25em}\textbf{（一）阅读 I（19 分）}\par
阅读下面的文字，完成 1～5 题。\par
\textbf{材料一：}\par
\hspace*{1.8em}“今天海况很不好，大家都在晕船。得知入选《自然》年度人物，我还来不及欢喜，因为更担心队员们的身体。\par
”电话里，传来杜梦然清脆的声音，“这个荣誉不是给我个人的，而是对中国深渊①科考团队的一份肯定。”\par
\hspace*{1.8em}2025 年 12 月 9 日，《自然》杂志公布 2025 年度十大人物，38 岁的中国科学院深海科学与工程研究所研究员杜梦\par
然入选。2024 年，她和团队在 9000 多米的深渊海底，发现了化能动物群，这是地球上已知最深的动物生态系统。此\par
后，通过更多航次的深潜科考，他们提出海底可能存在一个全球性的“化能生物走廊”。相关研究成果以论文与简报\par
形式发表于《自然》。\par
\hspace*{1.8em}杜梦然说，2024 年的那次深潜经历让她打破了自己的固有思维，也感受到生命的韧性。她从读博士时就开始研究\par
大陆架冷泉。“过去，一潜到海底，我就忙着找‘烟囱’，那是热液、冷泉、泥火山在海底‘冒泡’时制造的‘显眼\par
包’。但那次完全不同。”她回忆，那次下潜作业结束后，他们想从沟底爬上坡去看看，没想到发现了一整片化能动\par
物群——它们虽然看起来像植物，体内却共生着能够利用甲烷、硫化氢等化学物质获取能量从而维持自身生命活动的\par
微生物。“我们用深潜器碰了碰地面，十分坚硬，应该是冷泉活动次生的碳酸盐。”杜梦然说，原来深渊底部的冷泉\par
完全不同于大陆架冷泉，在千倍于海面的压力下，地底的甲烷、硫化氢等不可能以气体形式泄漏，只能以液体或水合\par
物的状态存在——这里的冷泉根本不冒泡。看到如此大规模的化能动物群，杜梦然不禁对大自然感到敬畏：“我们原\par
本以为，在这样的极端环境下，肯定不会有生命存在，但对化能生物来说，这里才是它们的乐园。”\par
\hspace*{1.8em}曾经，中国科学家要出海，不得不申请国外的航次，无法把采集的样品带回自己的实验室。这几年，杜梦然感到\par
，越来越多的外国科学家开始“追着我们跑”。“我们一个潜次带回的高质量样本，他们可能花一个星期、一个月都\par
拿不到。”杜梦然说，中国深渊科考团队用实力向全球证明，“我们可以创造一个又一个‘不可能’”。\par
\hspace*{1.8em}2025 年初，杜梦然和团队来到新西兰普伊斯哥海沟。那里地处“魔鬼西风带”，常年有六七个超强气旋活动，从\par
来没有人想过在这里下潜。“但我们去了，而且一口气下潜了 32 次。”利用风浪间歇的那短短几分钟，他们高效完\par
成了潜水器的布放和回收，创造了 75 小时 5 潜次的中国载人深潜新纪录。\par
\hspace*{1.8em}2025 年 6 月，由中国科学院深海科学与工程研究所牵头的国际大科学计划“全球深渊探索计划”正式获得联合国\par
“海洋科学促进可持续发展十年”执委会批准。杜梦然说，这标志着中国深渊科考开启全球合作新篇章，“我国的深\par
渊科考，在经历了多年‘跟跑’后，开始转向‘领跑’”。\par
\hspace*{1.8em}“科学面前，人人平等”，尽管投身深渊科学的女性很少，杜梦然却从不因为自己是女性就觉得需要被照顾，她\par
对自己的工作充满热爱。最让杜梦然开心的是，深渊科考是一个科学与工程紧密结合的领域，他们在船上经常要根据\par
作业需求，拆东补西地自己组装工具。“我们深海所有句话叫‘宁冒风险，不当逃兵’，研发出来的设备一定要在使\par
\clearpage
用中暴露不足，迭代改进，绝对不能摆进橱窗当摆设。”\par
\hspace*{1.8em}提出“全球性化能生物走廊”的假说，只是迈出了认识这个全新深渊生态系统的第一步。这次，杜梦然和团队要\par
下探的是一条从未去过的海沟。“虽然每个航次各有艰难，但每次都会让我们更有信心——什么大风大浪都经过了！\par
”\par
\hspace*{1.8em}（摘编自许琦敏《探寻深海未知，就该大胆潜下去》，《文汇报》2025 年 12 月 15 日）\par
\vspace{0.35em}\par
\textbf{材料二：}\par
\hspace*{1.8em}2024 年 7 月 8 日至 8 月 17 日，团队依托“探索 1 号”科考船与可下潜近 11000 米的全海深载人潜水器“奋斗者”\par
号，对千岛—堪察加海沟与西阿留申海沟沟底开展调查，在两条海沟均发现广泛分布的冷泉化能合成群落。在千岛—\par
堪察加海沟 FDZ271 潜次中，研究团队于 9533 米水深处首次发现密集化能合成群落，以管状蠕虫为主，栖息于沟底与\par
增生楔②底部交界处的黑色沉积物之上。\par
\hspace*{1.8em}在此次发现后，团队在相似地质环境开展了 23 潜次系统调查，以确定海沟底部化能合成群落的空间分布、规模\par
与生物多样性。在其中 19 次下潜中，研究人员对化能合成群落进行了观察、记录与采样，发现它们在沿着增生楔底\par
部的特殊区域——该区域沿着海沟底部延伸达 2500 千米——大量繁殖，构成了这两条深渊海沟的显著生态特征。尽\par
管前人在这些区域开展过沉积物和动物群调查，但该生态特征从未被发现和报道过。\par
\hspace*{1.8em}（摘译自彭晓彤、杜梦然等《深渊海沟最深处发现繁盛的化能合成生物》，《自然》2025 年 9 月 18 日，第 645 卷）\par
\vspace{0.35em}\par
[注] ①深渊：海洋深度大于 6000 米的区域。②增生楔：一种楔形地质体。\par
\vspace{0.35em}\par
1．下列对材料相关内容的理解和分析，不正确的一项是（3 分）（ ）\par
A．一次深潜作业后，杜梦然意外发现了深渊海底的化能动物群，也见识到深渊底部冷泉与大陆架冷泉的不同。\par
B．深渊底部的极端环境中依然可以有化能动物存在，这些动物体内共生着能够从化学物质中获取生命能量的微生物。\par
C．在杜梦然的科考团队下潜调查之前，人类未曾探索过千岛—堪察加海沟和地处“魔鬼西风带”的普伊斯哥海沟。\par
D．科考团队在相似地质环境的多次深潜调查，证实了海沟底部的化能合成生物群落并非零散分布，而是大规模存在着。\par
\vspace{0.35em}\par
2．根据材料内容，下列说法正确的一项是（3 分）（ ）\par
A．杜梦然说她感受到生命的“韧性”，这里的“韧性”既照应化能生物在极端环境下的存在，也表现了科学家坚持不\par
\hspace*{1.8em}懈、挑战极限的情怀。\par
B．在杜梦然的团队取得突破前，中国科学家的深渊科考只能依赖国外的航次，无法自主安排科研计划，采集的样品也\par
\hspace*{1.8em}无法带回。\par
C．研究团队在 9533 米水深处首次发现了密集的化能合成群落，根据这一发现，杜梦然团队提出了“全球性化能生物走\par
\hspace*{1.8em}廊”的假说。\par
D．材料二列出翔实的科考数据，例如“FDZ271 潜次”“23 潜次系统调查”“19 次下潜”“2500 千米”等，体现了严\par
谨的科学态度。\par
\vspace{0.35em}\par
3．材料一是一篇新闻作品，下列对其写作特点的分析，最恰当的一项是（3 分）（ ）\par
A．从杜梦然入选“《自然》年度十大人物”切入，聚焦于她的言行事迹，尤其注意引用她的话，写出了她对深渊探索\par
\hspace*{1.8em}事业的热爱。\par
B．第一段从“晕船”讲起，包含了何时、何地、何人、何故等新闻要素，有很强的现场感，也让读者了解到深渊科考\par
\hspace*{1.8em}的艰辛。\par
C．文学性较强，并不追求展示新闻事件的来龙去脉，而是截取 2024 年那次令杜梦然难忘的深潜经历，给予集中刻画。\par
D．以杜梦然团队的科研成就为例，辩证地讨论了科学精神的内涵，显现新闻时效性的同时，又有议论性文章的特点。\par
\vspace{0.35em}\par
4．中国深渊科考已开始转向“领跑”，这体现在哪些方面？请根据材料概括。（4 分）\par
\noindent 答：\rule{0.88\linewidth}{0.35pt}\par
\noindent 答：\rule{0.88\linewidth}{0.35pt}\par
5．某科技馆志愿者获得采访中国深渊科考团队的机会，采访话题围绕“全球深渊探索计划”展开，志愿者提出的\par
第一个问题是“计划的缘起与宗旨是什么”。请根据材料，再提出两个问题，并分别阐明提问的理由。（6 分）\par
\noindent 答：\rule{0.88\linewidth}{0.35pt}\par
\noindent 答：\rule{0.88\linewidth}{0.35pt}\par
\noindent 答：\rule{0.88\linewidth}{0.35pt}\par
\clearpage
\vspace{0.25em}\textbf{（二）阅读 II（20 分）}\par
阅读下面的文字，完成 6～9 题。\par
\vspace{0.35em}\par
\hspace*{1.8em}割草\par
\hspace*{1.8em}[俄]列夫   托尔斯泰\par
\hspace*{1.8em}列文是这样喜欢割草工作，今年春初以来，他就计划着整天和农民们一道去割草。他哥哥谢尔盖   伊万诺维奇到\par
来以后，他就踌躇起来。整天丢下哥哥一个人，他于心不安，又怕哥哥会为这事取笑他。但是当他走过草场，回想起\par
割草的印象，他几乎就决定要去割草了。在和哥哥激烈辩论之后，他又想到这个主意。\par
\hspace*{1.8em}“我需要体力活动，要不然，我的性情一定会变坏了。”他想，于是他下定决心去割草，不管在他哥哥或是农民\par
面前他会感到多么局促不安。\par
\hspace*{1.8em}傍晚，列文走到账房，派人到各村去召集明天的割草人手。\par
\hspace*{1.8em}“请把我的镰刀拿给季特，叫他磨好了明天给我，我也许要亲自去割草哩。”他说，竭力装得很安详的样子。\par
\hspace*{1.8em}①管家微微一笑，说：“好的，老爷。”\par
\hspace*{1.8em}晚上喝茶的时候列文对他哥哥说：“我看天气好起来了，明天我要开始割草了。”\par
\hspace*{1.8em}②“我很喜欢这种田间劳动。”谢尔盖   伊万诺维奇说。\par
\hspace*{1.8em}“我非常喜欢。明天我想要割一整天。”\par
\hspace*{1.8em}谢尔盖 伊万诺维奇抬起头来，“你是什么意思？像农民一样，从早到晚吗？”\par
\hspace*{1.8em}“是的，这是很愉快的。”列文说。\par
\hspace*{1.8em}“可是，农民们怎样看呢？他们一定会笑他们的主人是个怪物吧。”\par
\hspace*{1.8em}“不，我不这样想；人们无暇想到这些。”\par
\hspace*{1.8em}“但是你吃午饭怎么办呢？把你的红葡萄酒和烤火鸡送到那里未免有点尴尬吧。”\par
\hspace*{1.8em}“中午休息时我回来一趟就是了。”\par
\hspace*{1.8em}第二天早晨列文起得比平常早，但是他为了安排农场上的事耽搁了一会儿，当他到草场的时候，割草人已经在割\par
第二排了。\par
\hspace*{1.8em}他们在草场上高低不平的低处慢慢地刈割，列文认出了几个。这里，穿着白色长衬衫的叶尔米尔老头弯着腰在挥\par
着镰刀；那里，曾经做过列文马车夫的年轻小伙子瓦西卡把一排排的草一扫而光。还有季特，列文割草的师傅，一个\par
瘦小的农民。他在顶前面，大刀阔斧地割着，好像是在舞弄着镰刀一样。\par
\hspace*{1.8em}列文把马系在路旁，走到季特面前，季特从灌木丛里取出第二把镰刀，递给他。\par
\hspace*{1.8em}“磨好了，老爷；它像剃刀一样，自己会割哩。”季特说，带着微笑脱下帽子，把镰刀交给他。\par
\hspace*{1.8em}列文接了镰刀，试了试。割草的人们割完了一排，流着汗，愉快地、一个跟一个地走到路上来，微笑着和主人打\par
招呼。③他们都盯着他，但是没有人开口，直到一个身穿羊毛短衫的老头向他说话。\par
\hspace*{1.8em}“当心，老爷，可不要掉队啊！”他说，列文听到割草的人们中间压抑住的笑声。\par
\hspace*{1.8em}“我竭力不掉队就是了。”他说，站在季特背后，等待着开始割的时间。\par
\hspace*{1.8em}季特让出地方，列文就在他背后开始了。路边的草是短而坚韧的，列文被那么多眼睛注视着，弄得很狼狈，虽然\par
他使劲挥动着镰刀。他听到背后议论的声音：\par
\hspace*{1.8em}“镰刀把太高了，你看他的腰弯成那样。”有人说。\par
\hspace*{1.8em}“不要紧，他会顺手的，”老头子说，“他开了头了……”\par
\hspace*{1.8em}草渐渐柔软了，听着他们的话，列文没有回答，跟着季特，尽力割得好一点。他们前进了一百步。季特没有停步\par
，也没有露出丝毫疲惫的样子；但是列文已经开始担心自己要支持不下去了，他是这样地疲倦。\par
\hspace*{1.8em}他一面挥动着镰刀，一面感觉到他的气力已经使尽了。正在这时，季特停下了，弯下腰拾起一把草，擦净他的镰\par
刀，开始磨刀。列文伸直了腰，深深地舒了一口气，向四周望了一眼。季特磨快了自己的和列文的镰刀，他们又继续\par
前进。\par
\hspace*{1.8em}第二次还是一样。季特连续挥着镰刀没有停过，列文跟着他，竭力想不落在后面，他感觉到所有力气都用尽了，\par
④但是正在这个时候，季特又停下来磨镰刀。\par
\clearpage
\hspace*{1.8em}就这样他们割完了第一排。列文在他刈割的那块地面上走回来，这时候，尽管汗流满面，从鼻子上滴下，他的脊\par
背湿透得好像浸在水里一样，他还是感到非常愉快。特别使他高兴的是现在他知道他支持得了。\par
\hspace*{1.8em}只有一件事使他扫兴，就是他那一排割得不好。“我要少动胳膊，多用整个身子。”他想，拿季特那看去像切齐\par
了的一排，和自己那满地是草、参差不齐的一排比较着。\par
\hspace*{1.8em}往后几排就容易些了，但是列文还得使出全部力量才不至于落在农民后面。他什么也不想，耳朵里只听见镰刀的\par
飕飕声，眼前只看见季特渐渐远去的挺直的姿态，刈割了草的一片半圆形草地，在镰刀前面慢慢地像波浪一样倒下的\par
青草和花穗，以及前面可以休息的终点。\par
\hspace*{1.8em}突然，他感到他的热汗淋漓的肩膊上有一种愉快的凉爽感觉。他仰望了一下天空。阴沉的、低垂的乌云密布了，\par
大颗的雨点落下来。有的农民走去拿上衣穿上；有的农民，正如列文自己一样，只耸耸肩，享受着愉快的凉意。\par
\hspace*{1.8em}割完一排，又割一排。列文完全失去了时间观念，此刻天色是早是晚完全不知道了。他的工作开始发生了一种使\par
他非常高兴的变化。在劳动中竟有这样的时刻，他忘记了他在做什么，一切他都觉得轻松自如了，在这样的时候，他\par
那一排就割得差不多和季特的一样整齐出色了。但是他一想到他在做什么，而且开始竭力要做得好一些，他就立刻感\par
觉到劳动很吃力，而那一排也就割得不好了。\par
\hspace*{1.8em}又割了一排的时候，季特停下了，走到那老头跟前，低声对他说了句什么。两人都望了望天。“他们在谈什么呢\par
，为什么他们不接着割下去？”列文想，没有想到农民们已经刈割了四个多钟头没有休息，现在是他们吃早饭的时候\par
了。\par
\hspace*{1.8em}“吃早饭的时候了，老爷。”那老头子说。\par
\hspace*{1.8em}“已经是时候了吗？好的，那么吃早饭吧。”\par
\hspace*{1.8em}列文把镰刀交给季特，就和正要到放上衣的地方去拿面包的农民们一道，穿过一片被雨微微淋湿的刈割了的草地\par
，向他的马走去。这时他才想到，雨淋湿了他的干草。\par
\hspace*{1.8em}“干草会给糟蹋掉呢。”他说。\par
\hspace*{1.8em}“不会的，老爷；雨天割草晴天收嘛！”那老头子说。\par
\hspace*{1.8em}列文解下马缰，骑马回家去喝咖啡。\par
\hspace*{1.8em}谢尔盖 伊万诺维奇刚刚起来。列文喝完咖啡又回草场去了，而谢尔盖                                  伊万诺维奇还没有来得及穿好衣服走进\par
餐室。\par
\hspace*{1.8em}（周扬、谢素台译，有删改）\par
\vspace{0.35em}\par
[注] 本文节选自长篇小说《安娜           卡列宁娜》。文中列文是拥有田产、思想进步的贵族，他哥哥谢尔盖                        伊万诺维奇是位学者。\par
\vspace{0.35em}\par
6．下列对文本相关内容和艺术特色的分析鉴赏，不正确的一项是（3 分）（ ）\par
A．文本开头部分描写列文决心去割草的复杂心理过程，凸显列文亲自割草这件事在他所处的环境中并不寻常。\par
B．文本既细致地写出了列文的心理活动，同时又呈现了与列文生活相关的社会图景，气象开阔。\par
C．结尾先写列文和哥哥在同一个早上迥然不同的生活情形，这一细节呈现出兄弟二人不同的思想倾向。\par
D．列文奋力而忘我地割草，这一情节显示出他与《复活》中的聂赫留朵夫一样，都对底层人民怀有深深的忏悔感。\par
\vspace{0.35em}\par
7．关于文中画横线的语句，下列说法不正确的一项是（3 分）（ ）\par
A．句①，管家“微微一笑”，说明他对列文的含糊其辞和故作安详有自己的领会。\par
B．句②，谢尔盖         伊万诺维奇的话表明，他把田间劳动看作一种田园生活的趣味。\par
C．句③，农民们“没有人开口”，以此表达他们对主人前来割草一事无言的抗议。\par
D．句④，季特适时磨镰刀，说明他一直留意列文的进度，照顾列文的面子和体力。\par
\vspace{0.35em}\par
8．文中画波浪线的段落描写列文割草，读来令人感同身受。请结合文本谈谈这种真切感是如何表达出来的。（6\par
分）\par
\noindent 答：\rule{0.88\linewidth}{0.35pt}\par
\noindent 答：\rule{0.88\linewidth}{0.35pt}\par
\noindent 答：\rule{0.88\linewidth}{0.35pt}\par
9．原著中，列文早饭后回到草场继续割草，直到午饭时辰。下面就是吃午饭的场景描写，请结合对文本的阅读理\par
解，在括号处补写出一段文字，要求写出列文的举动和心理，衔接自然，字数在 60～100 字之间。（8 分）\par
\hspace*{1.8em}农民们集合了——从远处来的聚在大车下面，近的聚在铺着草的柳树下面。\par
\clearpage
\hspace*{1.8em}列文在他们旁边坐下；他不想走开了。\par
\hspace*{1.8em}在主人面前感到拘束的心情早已消失了。农民们预备午餐。老头子把一片面包捏碎，放进碗里，用匙柄捣烂，从\par
罐子里倒些水在上面，再捏一些面包进去，撒上一点盐。\par
\hspace*{1.8em}“哦，老爷，尝尝我的面包渣汤吧。”\par
\hspace*{1.8em}这面包渣汤是这么甘美，竟使列文放弃了回家吃饭的念头。他和老头子一道吃着，（ ）\par
\noindent 答：\rule{0.88\linewidth}{0.35pt}\par
\noindent 答：\rule{0.88\linewidth}{0.35pt}\par
\noindent 答：\rule{0.88\linewidth}{0.35pt}\par
\noindent 答：\rule{0.88\linewidth}{0.35pt}\par
\vspace{0.35em}\par
\vspace{0.25em}\textbf{（三）阅读 III（20 分）}\par
阅读下面的文言文，完成 10～14 题。\par
\textbf{材料一：}\par
\hspace*{1.8em}传曰：“子产治郑，民不能欺；子贱治单父，民不忍欺；西门豹治邺，民不敢欺。”三子之才能谁最贤哉？辨治\par
者当能别之。\par
\hspace*{1.8em}（节选自《史记        滑稽列传》）\par
\vspace{0.35em}\par
\textbf{材料二：}\par
\hspace*{1.8em}郑人游于乡校①，以论执政。然明②谓子产曰：“毁乡校，何如？”子产曰：“何为？夫人朝夕退而游焉，以议\par
执政之善否。其所善者，吾则行之；其所恶者，吾则改之。是吾师也，若之何毁之？我闻忠善以损怨，不闻作威以防\par
怨。岂不遽止？然犹防川，大决所犯，伤人必多，吾不克救也。不如小决使道，不如吾闻而药之也。”然明曰：“蔑\par
也今而后知吾子之信可事也。小人实不才若果行此其郑国实赖之岂唯二三臣？”\par
\hspace*{1.8em}（节选自《左传        襄公三十一年》）\par
\vspace{0.35em}\par
材料三：\par
\hspace*{1.8em}宓子贱治{\CJKfontspec{SimSun}亶父}③。三年，巫马旗短褐衣弊裘而往观化于{\CJKfontspec{SimSun}亶父}，见夜渔者，得则舍之。巫马旗问焉，曰：“渔为得\par
也，今子得而舍之，何也？”对曰：“宓子不欲人之取小鱼也。所舍者小鱼也。”巫马旗归，告孔子曰：“宓子之德\par
至矣，使民暗行若有严刑于旁。敢问宓子何以至于此？”孔子曰：“丘尝与之言曰：‘诚乎此者刑乎彼。’宓子必行\par
此术于{\CJKfontspec{SimSun}亶父}也。”\par
\hspace*{1.8em}（节选自《吕氏春秋          具备》）\par
\vspace{0.35em}\par
材料四：\par
\hspace*{1.8em}西门豹治邺，廪无积粟，府无储钱，库无甲兵，官无计会，人数言其过于文侯。文侯身行其县，果若人言。文侯\par
曰：“翟璜任子治邺而大乱。子能道，则可；不能，将加诛于子。”西门豹曰：“臣闻王主富民，霸主富武，亡国富\par
库。今君欲为霸王者也，臣故蓄积于民。君以为不然，臣请升城鼓之，甲兵粟米可立具也。”于是乃升城而鼓之。一\par
鼓，民被甲括矢，操兵弩而出；再鼓，负辇粟而至。文侯曰：“罢之。”西门豹曰：“与民约信，非一日之积也。一\par
举而欺之，后不可复用也。燕常侵魏八城，臣请北击之，以复侵地。”遂举兵击燕，复地而后反。\par
\hspace*{1.8em}（节选自《淮南子         人间训》）\par
\vspace{0.35em}\par
[注] ①乡校：古代地方学校。②然明：郑大夫{\CJKfontspec{SimSun}鬷}蔑的字。③亶父：即单父，春秋时鲁邑。\par
\vspace{0.35em}\par
10．材料二画波浪线的部分有三处需要断句，请用铅笔将答题卡上相应位置的答案标号涂黑，每涂对一处给 1 分，\par
涂黑超过三处不给分。（3 分）\par
小人 A 实不才 B 若果 C 行 D 此 E 其郑国 F 实 G 赖之 H 岂唯二三臣？\par
11．下列对材料中加点的词语及相关内容的解说，不正确的一项是（3 分）（ ）\par
A．善，褒扬，与《鸿门宴》中“素善留侯张良”的“善”用法、意思均相同。\par
B．药，以……为药，与《赤壁赋》中“侣鱼虾而友麋鹿”的“侣”“友”用法相同。\par
C．身，亲身、亲自，与《陈涉世家》中“将军身被坚执锐”的“身”意思相同。\par
D．富，使……富足，与《归去来兮辞》中“眄庭柯以怡颜”的“怡”用法相同。\par
\clearpage
12．下列对材料有关内容的概述，不正确的一项是（3 分）（ ）\par
A．郑国人常常聚在乡校里议论执政，大夫然明建议毁掉乡校，子产认为不仅不能毁掉乡校，而且应把那些议论者作为\par
\hspace*{1.8em}自己的老师，然明对子产的见解十分佩服。\par
B．巫马期感慨于宓子贱治理亶父的成效，并就其原因向孔子请教。孔子表示自己告诫过宓子贱，宓子贱定是按照以诚\par
\hspace*{1.8em}化民的主张去治理亶父的。\par
C．西门豹担任邺令时，粮仓没有粮食，府库没有储钱，武库没有兵器，官衙没有账目，有人屡次向魏文侯告状。文侯\par
\hspace*{1.8em}巡视后发现情况果真如此。\par
D．西门豹向文侯表示只要自己登城击鼓，就能马上备好武器粮食，两通鼓果然聚其所需。文侯见此，免去了西门豹的\par
邺令，改派他领兵去攻打燕国。\par
\vspace{0.35em}\par
13．把材料中画横线的句子翻译成现代汉语。（8 分）\par
（1）渔为得也，今子得而舍之，何也？\par
\noindent 答：\rule{0.88\linewidth}{0.35pt}\par
\noindent 答：\rule{0.88\linewidth}{0.35pt}\par
（2）燕常侵魏八城，臣请北击之，以复侵地。\par
\noindent 答：\rule{0.88\linewidth}{0.35pt}\par
\noindent 答：\rule{0.88\linewidth}{0.35pt}\par
14．材料一中说“西门豹治邺，民不敢欺”，你认为材料四能否印证这一说法？请说明理由。（3 分）\par
\noindent 答：\rule{0.88\linewidth}{0.35pt}\par
\noindent 答：\rule{0.88\linewidth}{0.35pt}\par
\vspace{0.35em}\par
\vspace{0.25em}\textbf{（四）阅读 IV（15 分）}\par
阅读下面这首宋诗，完成 15～17 题。\par
\vspace{0.35em}\par
\hspace*{1.8em}风\par
\hspace*{1.8em}司马光\par
\vspace{0.35em}\par
\hspace*{1.8em}陇首红旌急，樽前翠幕重。\par
\hspace*{1.8em}林端翻远刹，花外转疏钟。\par
\hspace*{1.8em}夜寂清机[注]发，春阑别意浓。\par
\hspace*{1.8em}如何玉琴韵，并欲在青松。\par
[注] 清机：清净的心思。\par
\vspace{0.35em}\par
15．下列对这首诗的理解和赏析，不正确的一项是（3 分）（ ）\par
A．诗中妙用“翻”字，表现远处的庙宇因林梢摇动，若隐若现。\par
B．花丛外随风传来的钟声依稀缥缈，用“疏钟”二字尤为贴切。\par
C．琴曲中传达出松涛的韵味，缘自人的内心受到了风的触动。\par
D．诗人通过对风的描写，表达出远离尘嚣、归隐山林的心愿。\par
\vspace{0.35em}\par
16．有一首歌这样唱道：“如果你要写风，就不能只写风/你要写柳条轻轻柔柔飘入你心中/竹筏轻轻飘荡，湖面在\par
泛粼光/船上悬挂的铃儿响呀响呀响叮当……你要用眼睛耳朵捕捉每一刻/用心感受。”请具体阐述这首歌与本诗的\par
相通之处。（6 分）\par
\noindent 答：\rule{0.88\linewidth}{0.35pt}\par
\noindent 答：\rule{0.88\linewidth}{0.35pt}\par
\noindent 答：\rule{0.88\linewidth}{0.35pt}\par
17．补写出下列句子中的空缺部分。（6 分）\par
（1）本诗中的“翠幕”指青绿色的帷幕，柳永《望海潮》中“\rule{0.080\linewidth}{0.35pt}，\rule{0.080\linewidth}{0.35pt}”两句写街巷\par
河桥的美丽以及房舍宅院的雅致，也用到了这个词语。\par
\clearpage
（2）在古代诗歌中，风常常与雨、雪等对举，如陆游《书愤》中的“\rule{0.080\linewidth}{0.35pt}，\rule{0.080\linewidth}{0.35pt}”。\par
（3）在古代散文中，写到风的名句也是层出不穷，比如：“\rule{0.080\linewidth}{0.35pt}，\rule{0.080\linewidth}{0.35pt}。”\par
\vspace{0.35em}\par
\vspace{0.35em}\par
\vspace{0.25em}\textbf{二、语言文字运用（16 分）}\par
阅读下面的文字，完成 18～21 题。\par
\hspace*{1.8em}“手机随手一拍，竞然有点像电影剧照。”不少人刷到朋友圈或社交媒体平台上的“人生壁纸”，都会忍不住多\par
看两眼。有人在幕色中的地铁站前仰望天空，有人在雨夜的路灯下形色匆匆，也有人在洒满阳光的院子里赏花喝茶\par
……一组媒体调查数据显示，89.5\%的受访者有记录生活的习惯，其中超半数会通过拍照留存生活瞬间。社交媒体平\par
台上，怎么拍才能“出片”、哪些地方最“出片”这类话题也持续火曝。\par
\hspace*{1.8em}有人点赞：“喜欢认真生活的样子。”也有人疑惑：“只为拍张照片，至于吗？”对“出片”的不同理解，折射\par
出人们对生活的多元探寻。\par
    事实上，“出片”不是为了展示精致的构图、完美的滤镜，而是源于对生活的热爱。它\rule{0.080\linewidth}{0.35pt}，\rule{0.080\linewidth}{0.35pt}，\rule{0.080\linewidth}{0.35pt}\par
，\rule{0.080\linewidth}{0.35pt}。\par
\hspace*{1.8em}这份记录生活的热情，不仅丰富了个体的精神世界，有时还在现实中创造出意想不到的价值。广东广州，大学生\par
随手拍摄的街景，在海外社交平台收获超百万浏览量，向世界展现了中国城市的美好与烟火气；浙江永嘉，摄影爱好\par
者分享的柿子林照片意外走红，当地顺势打造出一个很“出片”的千亩柿子林休闲景区，带动村集体及村民增收约\par
150 万元。这些定格在影像中的美好，不经意间成为展现城市风貌、助力乡村振兴的鲜活见证。\par
\hspace*{1.8em}然而，如果“出片”从纯粹的生活记录和发自内心的情感表达，异化为某种刻意追求，便容易偏离初衷、变了味\par
道。①为了“出片”，有人不顾危险闯入禁区；②有人霸占打卡点妨碍交通；③甚至因盲目追求打卡而引发安全事故\par
也时有发生。④当“出片”的执念突破了安全与文明的底线，⑤因此失去了记录生活的意义。\par
18．文中第一段有四个含错别字的词语，请找出并改正。（4 分）\par
\noindent 答：\rule{0.88\linewidth}{0.35pt}\par
\noindent 答：\rule{0.88\linewidth}{0.35pt}\par
19．依次填入文中横线处的语句，衔接最恰当的一组是（3 分）（ ）\par
①承载着对生活仪式感的珍视\par
②看似寻常的瞬间被定格成画面\par
③没有专业门槛，不苛求拍摄技巧\par
④更蕴含着人们热爱生活的真情实感\par
⑤是普通人用镜头发现美好、用影像寄托情感的自然表达\par
A．①⑤③④② B．③⑤②①④ C．⑤③①④② D．②①⑤③④\par
\vspace{0.35em}\par
20．文中画波浪线的部分有两处表述不当，请指出其序号并修改，使语言准确流畅，逻辑严密，可少量增删词语，\par
不得改变原意。（4 分）\par
\noindent 答：\rule{0.88\linewidth}{0.35pt}\par
\noindent 答：\rule{0.88\linewidth}{0.35pt}\par
21．《现代汉语大词典》已经收录了“朋友圈”“打卡”等反映时代特点的新词。“出片”作为近年进入大众生活\par
的热词，如果也要收入词典，你会如何释义？要求写出两个义项，并给每个义项拟一个相应的例句，语言准确简洁。\par
（5 分）\par
\noindent 答：\rule{0.88\linewidth}{0.35pt}\par
\noindent 答：\rule{0.88\linewidth}{0.35pt}\par
\noindent 答：\rule{0.88\linewidth}{0.35pt}\par
\noindent 答：\rule{0.88\linewidth}{0.35pt}\par
\clearpage
\vspace{0.25em}\textbf{三、写作（60 分）}\par
22．阅读下面的材料，根据要求写作。（60 分）\par
\hspace*{1.8em}词语是表达思想情感的载体，也是展现社会生活变化的窗口。当前，世界之变、时代之变、历史之变正以前所未\par
有的方式展开。青年是常为新的，在你的成长过程中，你对哪一个词语的理解发生了变化？这变化有你成长的印记，\par
对你有特殊的意义……\par
\hspace*{1.8em}以上材料引发了你怎样的联想和思考？请写一篇文章。\par
\hspace*{1.8em}要求：选准角度，确定立意，明确文体，自拟标题；不要套作，不得抄袭；不得泄露个人信息；不少于 800 字。\par
\vspace{0.35em}\par
\vspace{0.35em}\par
整理校核依据：教育部教育考试院发布的 2026 年高考语文试题评析及公开创新试题材料；中国全民阅读网转载的命题组专家解析；公\par
开完整试卷文字转录。\par
\clearpage
\end{document}
